\section{Introducción}
El reparto del caudal en una bifurcación de canales abiertos determina el dimensionamiento de canales de riego, derivaciones y obras de control. La literatura experimental y numérica muestra que la fracción derivada depende del ángulo, de la relación de anchos, del número de Froude y de las condiciones aguas abajo, y que las zonas de separación y recirculación junto a la esquina de la derivación modifican la sección efectiva de flujo \citep{ramamurthy1988,ramamurthy1990,hsu2002,momplot2017,ibrahim2022,gumgum2026}. Los modelos unidimensionales (1D) resuelven la unión mediante continuidad, igualdad de niveles y pérdidas locales, y son baratos y fáciles de integrar en redes \citep{chaudhry2008,chow1959}; los modelos bidimensionales (2D) capturan la redistribución lateral, pero dependen de la malla, de la rugosidad y de las condiciones de borde \citep{ghostine2012,ramamurthy2007}.

Para una bifurcación dada, el ingeniero necesita saber (i) qué formulación 1D reproduce el reparto con un error aceptable, (ii) de qué parámetros depende más ese reparto y (iii) si las diferencias frente a un modelo 2D son un defecto de la formulación o un efecto de configuraciones no equivalentes. Las comparaciones 1D--2D publicadas suelen reportar el acuerdo, pero no separan estas causas \citep{ghostine2012,gumgum2025}.

Este trabajo aborda tres objetivos para un caso sintético: (1) resolver por Newton--Raphson dos formulaciones del reparto y verificar el código; (2) cuantificar la sensibilidad del reparto a los parámetros y propagarla a una incertidumbre paramétrica; y (3) contrastar las dos formulaciones con simulaciones HEC-RAS 2D en cinco ángulos, expresando el sesgo con criterios de aceptación declarados antes de mirar los resultados y combinándolo con la sensibilidad para decidir si el sesgo es compatible con la incertidumbre de los parámetros. El aporte no es un modelo nuevo, sino un procedimiento reproducible para decidir cuándo una formulación 1D de la unión es suficiente y cuándo no, junto con una lista explícita de lo que la comparación efectuada no permite concluir.

El artículo sigue el orden habitual de Water: marco teórico (Sección~\ref{sec:teoria}), materiales y métodos (Sección~\ref{sec:metodos}), resultados (Sección~\ref{sec:resultados}), discusión (Sección~\ref{sec:discusion}) y conclusiones (Sección~\ref{sec:conclusiones}).
